\chapter{Conclusion}
\label{ch:conclusion}

\section{Answers to the Research Questions}

This thesis asked two questions: how a frequent batch auction engine is
designed, and what changes when the same real order flow is cleared in
batches instead of continuously. The two paragraphs below answer each in
turn.

\paragraph{RQ1.} A frequent batch auction engine is fully specified by five
design blocks: the batch interval, the disclosure regime during
accumulation, the order flow segregation architecture, the clearing and
rationing rule, and settlement and residual handling. These are not
optional extras layered onto the mechanism but decisions that serial
processing settles automatically in a continuous engine, and which become
visible and contestable only once it is abandoned. \Cref{tab:consequences}
lists every option with its predicted consequence. Three points carry the
most weight. The only unambiguous structural gain is discrete accumulation
combined with uniform pricing; every other choice trades one distortion for
another, so the question is which distortion the venue prefers. The blocks
interact and cannot be optimised independently. And a venue can implement
frequent batch auctions exactly as proposed and still have a speed race,
through two separate channels: price-time rationing inside the batch
rewards being early, and a fixed, predictable close rewards being late ---
the sharpest practical implication of the analysis.
Two questions with no settled answer were developed as contributions: the
maker--taker distinction dissolves under simultaneous uniform-price
clearing, with persistence across batches the strongest of four candidate
redefinitions; and the price improvement uniform clearing produces is a
transfer between buyers and sellers, not a gain, whose split is set by the
clearing-price selection rule. Under a fixed close and a visible batch,
that split can be steered by whoever submits last, which turns the choice
of time window into a boundary repricing race (\Cref{sec:priceimprovement}).

\paragraph{RQ2.} The simulation replayed one month of SOL perpetual order
flow through both engines, with one-second batches, and compared them
second by second on the same orders. It produced six findings
(\Cref{tab:findings}); \Cref{sec:findingsreliability} explains how each is
reached, why it arises and how far it can be relied on.
\begin{enumerate}[label=\textbf{F\arabic*}]
  \item \textbf{The same volume trades, at the same price} (high
  reliability). Batching reschedules trades rather than preventing them:
  unfilled orders roll over, so crossing interest still meets, at most one
  interval later. The two prices differ by a median of
  \PriceDiffMedianBps\,bps, and only when an order sweeps several price
  levels, which the CDA prices level by level and the FBA at one price.
  \item \textbf{The price responds to order flow by the same amount}
  (high reliability at 5\,s). Net buying removes the same resting orders in
  both engines, so both prices settle at the same level. Theory predicts a
  smaller response under batching through a change in behaviour, which a
  replay cannot reproduce.
  \item \textbf{Uniform pricing removes price dispersion within a second}
  (high reliability). Each one-second bucket holds one batch, which trades
  at one price, so all traders arriving within that second get the same
  price; the variation moves to between seconds.
  \item \textbf{At one second batching adds no liquidity} (high
  reliability for no gain, low for a loss). Depth near the price is not
  higher, but a few percent lower; Amihud illiquidity and volatility are
  slightly higher, partly because of how each engine's price is measured.
  A one-second batch rarely holds more than the flow that would have traded
  anyway, so there is little to pool.
  \item \textbf{Batching delays execution} (high reliability for the
  direction). Orders wait \ResTimeToExecDiff\,s longer for their first
  fill: up to one second for the clear, the rest most likely from which
  resting orders the batch matches. Most batches clear no trade, but the
  continuous market also trades in only about a quarter of seconds.
  \item \textbf{The continuous market charges resting orders an
  adverse-selection cost} (high reliability, continuous market only). The
  realized spread is negative and the price impact grows with the horizon:
  resting orders are accepted most often by traders who expect the price
  to move their way.
\end{enumerate}

In summary, at one second and on this market, batching largely
reproduces the continuous market with a short delay. Its promised gains
depend on pooling orders that would otherwise arrive at different
moments, and a one-second batch rarely holds more than the flow that would
have traded anyway. This agrees with field evidence from the Taiwan Stock
Exchange: \citet{lee2026taiwan} find that liquidity and price efficiency
improved for mid- and small-cap stocks when the exchange moved from
frequent batch auctions to continuous trading. Their design captures the
participants' response to the mechanism, this one the mechanism alone,
and neither finds a liquidity gain from batching. Whether batching
protects resting orders from adverse selection, and whether a longer
interval would give the batches enough flow to pool, remain open
(\Cref{sec:limitssign}).

\section{Limitations}
\label{sec:limitations}

\paragraph{Participants cannot respond.}
\label{sec:limitsbehaviour}
A replay holds order flow constant, which is what makes the comparison
clean, and by the same token prevents participants from responding to the
mechanism. The replayed orders were submitted to a continuous engine, so
their sizes, arrival times and limit prices were not chosen against a batch
rule. The equilibrium predictions of the FBA literature --- narrower quoted
spreads, size inflation, orders bunching at the batch close, bid shading
(\Cref{sec:priceimprovement}) --- are therefore out of reach, and the
findings are the mechanical baseline against which such responses would be
measured. The same holds for cancellations: if latency races were present,
a continuous market would show more cancellations than a batch market, but
the data contain only the continuous market's cancellations.

\paragraph{One venue, one instrument, one month, one interval.} Magnitudes
should not be carried over to other markets without re-estimation.
Hyperliquid processes cancellations before new orders within each block
(\Cref{sec:venue}), so its continuous market already protects resting
quotes to a degree; on a purely first-come, first-served exchange the
benefit of batching could be larger. And every finding is for
$\tau = 1$\,s.

\paragraph{Fidelity of the replay.} The replay departs from the venue in
ways that affect both engines alike --- unenforced post-only and
self-trade rules, rounding to whole SOL, and restarts in the first days.
None favours either side of the paired comparison, but absolute levels
describe the replay, not the live market (\Cref{sec:limitsreplay}).

\section{Future Work}
\label{sec:futurework}

Each limitation above points to a concrete extension.
\begin{itemize}
  \item \textbf{Measure adverse selection under batching.} Recording the
  signed deviation of each FBA trade from its reference price, and signing
  it by the venue's net taker flow in that second, would give the FBA a
  spread decomposition comparable with the CDA's (\Cref{sec:limitssign})
  and answer the question F6 leaves open.
  \item \textbf{Vary the interval.} A sweep over $\tau$ would show whether
  longer intervals fill the batches without losing F1 and F2, and the limit
  $\tau \to 0$ would check that the FBA converges to the CDA.
  \item \textbf{Replay the month without interruption}, which removes the
  effects of the restarts (\Cref{sec:testresets}).
  \item \textbf{Let participants respond.} Agent-based simulation or
  laboratory experiments built on the same engines could test the
  behavioural predictions a replay cannot reach.
  \item \textbf{Other instruments and venues}, to see whether the findings
  hold beyond one perpetual contract on one exchange.
\end{itemize}
